%%%%%%%%%%%%%%%%%%%%%%%%%%%%%%%%%%%%%%%%%%%%%%%%%%%%%%%%%%%%%%%%%%%%%%%%%%%%%%%%
% Preámbulo                                                                    %
%%%%%%%%%%%%%%%%%%%%%%%%%%%%%%%%%%%%%%%%%%%%%%%%%%%%%%%%%%%%%%%%%%%%%%%%%%%%%%%%

\documentclass[11pt,a4paper,titlepage,oneside]{report}

%%% RELACIÓN DE VARIABLES A PERSONALIZAR %%%
\def\lingua{gal}
%\def\lingua{esp} % descomenta esta liña se redactarás a memoria en español
%\def\lingua{eng} % descomenta esta liña se redactarás a memoria en inglés
\def\nome{Hugo Mato Cancela}                             % substitúe aquí o teu nome
\def\nomedirectorA{Roberto Rodríguez Osorio}              % substitúe aquí o nome de quen dirixe
%\def\nomedirectorB{Outro Nome Completo}             % duplica esta liña máis veces se o precisas, cambiando
                                                     % a letra final (A, B, C, D...): úsanse na portada.tex
\def\titulo{Simulador de RISC-V empregando SystemC} % substitúe aquí o título do teu TFG
%\def\titulacion{gced}                               % descomenta esta liña e comenta a seguinte se es estudante do GCED
\def\titulacion{gei}
%\def\mencion{COMPUTACIÓN}                           % descomenta a mención que che corresponda se es estudante do GEI
%\def\mencion{ENXEÑARÍA DO SOFTWARE}
\def\mencion{ENXEÑARÍA DE COMPUTADORES}
%\def\mencion{SISTEMAS DE INFORMACIÓN}
%\def\mencion{TECNOLOXÍAS DA INFORMACIÓN}

%\def\renomearcadros{si} % descomenta esta liña se redactas a memoria en español e prefires que
                         % os "cuadros" e o "índice de cuadros" se renomeen
                         % a "tablas" e "índice de tablas" respectivamente

\usepackage{estilo_tfg}

% Lista de paquetes potencialmente interesantes (uso baixo demanda)

% \usepackage{alltt}       % proporciona o entorno alltt, semellante a verbatim pero que respecta comandos
% \usepackage{enumitem}    % permite personalizar os entornos de lista
% \usepackage{eurofont}    % proporciona o comando \euro
% \usepackage{float}       % permite máis opcións para controlar obxectos flotantes (táboas, figuras)
% \usepackage{hhline}      % permite personalizar as liñas horizontais en arrays e táboas
  \usepackage{longtable}   % permite construir táboas que ocupan máis dunha páxina
% \usepackage{lscape}      % permite colocar partes do documento en orientación apaisada
% \usepackage{moreverb}    % permite personalizar o entorno verbatim
% \usepackage{multirow}    % permite crear celdas que ocupan varias filas da mesma táboa
% \usepackage{pdfpages}    % permite insertar ficheiros en PDF no documento
% \usepackage{rotating}    % permite diferentes tipos de rotacións para figuras e táboas
% \usepackage{subcaption}  % permite a inclusión de varias subfiguras nunha figura
% \usepackage{tabu}        % permite táboas flexibles
% \usepackage{tabularx}    % permite táboas con columnas de anchura determinada

%%%%%%%%%%%%%%%%%%%%%%%%%%%%%%%%%%%%%%%%%%%%%%%%%%%%%%%%%%%%%%%%%%%%%%%%%%%%%%%%
% Corpo                                                                        %
%%%%%%%%%%%%%%%%%%%%%%%%%%%%%%%%%%%%%%%%%%%%%%%%%%%%%%%%%%%%%%%%%%%%%%%%%%%%%%%%

\begin{document}

 %%%%%%%%%%%%%%%%%%%%%%%%%%%%%%%%%%%%%%%%
 % Preliminares do documento            %
 %%%%%%%%%%%%%%%%%%%%%%%%%%%%%%%%%%%%%%%%

 \begin{titlepage}
  
  \hspace*{128pt}
  \textcolor{udcpink}{{\fontencoding{T1}\fontfamily{phv}\selectfont Facultade de Informática}}\\[-32pt]

  \begin{center}
    \includegraphics[scale=0.3]{imaxes/udc}\\[25pt]

    {\large TRABALLO FIN DE GRAO \\
            \nometitulacion \\
            \nomemencion } \\[10pt]

    \carimbo \\[25pt]

    \begin{huge}
      \begin{spacing}{1.3}
        \bfseries \titulo
      \end{spacing}
    \end{huge}
  \end{center}
  
  \vfill
  
  \begin{flushright}
    {\large
    \begin{tabular}{ll}
      {\bf Estudante:} & \nome \\
      {\bf Dirección:} & \nomedirectorA \\
%                      & \nomedirectorB \\ % duplica esta liña máis veces se o precisas, cambiando
                                           % a letra final (A, B, C, D...); define eses nomes no memoria_tfg.tex
    \end{tabular}}
  \end{flushright}
  \rightline{A Coruña, \datasimple.}
\end{titlepage}

 \dedicatoria{Dedicatoria} % escribe neste comando o teu texto de dedicatoria
 \paxinaenbranco
 \begin{agradecementos}
 \blindtext                % substitúe este comando polo teu texto de agradecementos
 \end{agradecementos}
 %%%%%%%%%%%%%%%%%%%%%%%%%%%%%%%%%%%%%%%%%%%%%%%%%%%%%%%%%%%%%%%%%%%%%%%%%%%%%%%%

\pagestyle{empty}
\begin{abstract}
      RISC-V é unha nova arquitectura libre para procesadores programables. Está chamada a competir con outras arquitecturas máis establecidas como \acrfull{arm} en moitos ámbitos tecnolóxicos. Ademais de que as especificacións de RISC-V son abertas, a principal vantaxe desta arquitectura é que cubre desde as implementacións máis sinxelas para sistemas embarcados, ata as máis potentes para cálculo científico e multimedia. Todo isto é posible mediante a especificación de moitas das características máis avanzadas como extensións ás arquitecturas básicas. 

A implementación dun procesador require previamente dun modelado e simulación que garantan que o funcionamento final vai ser o correcto. 

Neste traballo, pártese dun modelo da versión básica RV32I realizada en SystemC. A única extensión implementada é a multiplicación para enteiros. O obxectivo deste proxecto é modelar e simular extensións adicionais ata chegar ao coñecido como nivel G. 

  \vspace*{25pt}
  \begin{segundoresumo}
    RISC-V is a new open-source architecture for programmable processors. It is destined to compete with more established architectures like \acrshort{arm} in a lot of technological fields. In addition to the fact that the RISC-V specifications are open, the main advantage of this architecture is that it spans from the simplest implementations for embedded systems to the most powerful ones for scientific computing and multimedia. All of this is possible through the specification of many of the most advanced features as extensions to the basic architectures.

The implementation of a processor requires a previous modeling and simulation that ensures that its final operation will be correct.

This work is based on a model of the basic RV32I implemented in SystemC. The only extension implemented is integer multiplication. The objective for this project is to model and simulate additional extensions up to what is known as the G level.
  \end{segundoresumo}
\vspace*{25pt}
\begin{multicols}{2}
\begin{description}
\item [\palabraschaveprincipal:] \mbox{} \\[-20pt]
  \begin{itemize}
      \item RISC-V
      \item Simulador
      \item SystemC
      \item Procesador
      \item Arquitectura
      \item Extensións
      \item C++
  \end{itemize}
\end{description}
\begin{description}
\item [\palabraschavesecundaria:] \mbox{} \\[-20pt]
  \blindlist{itemize}[7] % substitúe este comando por un itemize
                         % que relacione as palabras chave
                         % que mellor identifiquen o teu TFG
                         % no idioma secundario da memoria (tipicamente: inglés)
\end{description}
\end{multicols}

\end{abstract}
\pagestyle{fancy}

%%%%%%%%%%%%%%%%%%%%%%%%%%%%%%%%%%%%%%%%%%%%%%%%%%%%%%%%%%%%%%%%%%%%%%%%%%%%%%%%


 \pagenumbering{roman}
 \setcounter{page}{1}
 \bstctlcite{IEEEexample:BSTcontrol}

 \tableofcontents
 \listoffigures
 \listoftables
 \clearpage
 
 \pagenumbering{arabic}
 \setcounter{page}{1}

 %%%%%%%%%%%%%%%%%%%%%%%%%%%%%%%%%%%%%%%%
 % Capítulos                            %
 %%%%%%%%%%%%%%%%%%%%%%%%%%%%%%%%%%%%%%%%

 \chapter{Introdución}
\label{chap:introducion}

\lettrine{E}{ste} proxecto busca crear un simulador de RISC-V empregando a librería SystemC en C++. Ao longo desta memoria describiranse as etapas de execución, os distintos módulos e a motivación destes así como as extensións.


%%citas~\cite{ErlangBook,ErlangWebBook} e de referencias
%%internas (sección \ref{sec:mostra}, páxina \pageref{sec:mostra}).


\section{Motivación}
\label{sec:motivación}
RISC-V apunta a ser unha das arquitecturas máis empregadas nun futuro, xa que é libre, permitindo aforrar o custo de licenzas. Grazas a que se pode modificar, engadindo ou eliminando funcionalidades, isto permite que abarque múltiples sectores, dende chips máis sinxelos orientados a IoT, ata competir con ARM en sistemas embebidos. REFES AQUI\\
Durante o proceso de deseño, unha parte clave é a verificación do correcto funcionamento. REFES AQUI Se ben é posible crear un chip con cada versión, na práctica, debido aos longos tempos e altos custos, é inviable. Ahí é onde un simulador toma protagonismo, xa que permite probar de forma rápida, simple e barata os deseños creados. 

\section{Obxectivos}
\label{sec:obxectivos}
Os obxectivos deste proxecto son modelar e simular, usando SystemC, as seguintes extensións da arquitectura RV32I: 

\begin{itemize}
    \item multiplicación e división por números enteiros (extensión M).
    \item aritmética en punto flotante de simple (extensión F).
    \item  operacións atómicas  (extensión A).
    \item  xestión de rexistros de control e estado (extensión Zicsr).
    \item sincronización de escritura de instrucións (extensión Zifencei).
\end{itemize}

O modelado será totalmente parametrizable, permitindo especificar a latencia das diferentes instrucións. Tamén permitirá especificar o número de canles de execución para as unidades de enteiros a punto flotante, e se estes están ou non totalmente segmentados. 

Desta maneira, a simulación permitirá comparar o rendemento de, por exemplo, unha implementación na que multiplicador e divisor comparten circuítos, cunha na que ambos son independentes, e tamén comparar un divisor totalmente segmentado con un que non o sexa. 

Os resultados do modelado e a simulación son dous: verificar o correcto funcionamento da arquitectura, e comprobar o seu rendemento. 


\section{Metodoloxía}
\label{sec:metodoloxía}
O método de traballo será incremental, dividindo as tarefas en partes independentes que van ser implementadas, simuladas e verificadas por orde de complexidade antes de proceder coa seguinte. \\
O procedemento habitual é unha reunión semanal na que se revisa o feito na semana anterior, acompañado de probas cos tests correspondentes para esa parte. Despois, decídese cal é o seguinte paso, podendo ser a implementación dunha nova extensión ou modificar un módulo do simulador.

\subsection{Fases principais}
\begin{itemize}
    \item Estudio da documentación existente sobre RISC-V.
    \item Familiarización coa implementación base de RV32I en SystemC. 
    \item Modelado e simulación do multiplicador e divisor de enteiros. 
    \item Modelado e simulación das extensións de punto flotante F e D. 
    \item Modelado e simulación de extensións Zicsr e  Zifencei.
    \item Empaquetamento do software. 

\end{itemize}

\section{Contida da memoria}
\label{sec:contido_memoria}
Texto aquí :)


 \chapter{Contido demostrativo}
\label{chap:demo}

\lettrine{E}{ntre} a introdución e as conclusións, o documento conterá
tantos capítulos como sexa preciso, sempre con coidado de non rebasar
o límite de 80 páxinas fixado polo regulamento de TFGs.

Empregaremos éste de xeito demostrativo, para ilustrar o uso de
elementos habituais que poidan ser de utilidade\footnote{Por exemplo,
  isto é unha nota a pé de páxina.}.

\section{Inclusión de imaxes}

Se precisamos imaxes no noso documento, incluirémolas do xeito que se
indica na figura~\ref{fig:exemplo} (páxina~\pageref{fig:exemplo}). Se
o facemos así, \LaTeX ubicará cada imaxe no mellor lugar posible,
lugar que pode variar a medida que o documento vaia crecendo coa
inclusión de máis texto e outros elementos (máis imaxes, táboas,
etc.).

\begin{figure}[hp!]
  \centering
  \includegraphics[width=0.75\textwidth]{imaxes/udc.png}
  \caption{Pé de imaxe descritivo}
  \label{fig:exemplo}
\end{figure}

Recoméndase almacenar os ficheiros gráficos no directorio
\texttt{imaxes}.

\subsection{Inclusión de varias sub-imaxes}

Se precisamos inserir imaxes relacionadas, pode ser apropiado
incluílas como sub-figuras, do xeito que se pode apreciar na
figura~\ref{fig:exemplo-subfiguras} coas
imaxes~\ref{fig:subfigura-rotada}
e~\ref{fig:subfigura-deformada}. Como se pode ver nos exemplos desta
sección, sempre é recomendable referirse ás imaxes pola súa
referencia, xa que dese xeito non dependemos de onde queden ubicados
os elementos en cuestión.

\begin{figure}[hp!]
  \centering
  \begin{subfigure}[c]{0.3\textwidth}
    \includegraphics[angle=45,width=\textwidth]{imaxes/udc.png}
    \caption{Pé de subimaxe rotada}
    \label{fig:subfigura-rotada}
  \end{subfigure}
  \hspace{0.1\textwidth}
  \begin{subfigure}[c]{0.3\textwidth}
    \includegraphics[width=\textwidth,height=3cm]{imaxes/udc.png}
    \caption{Pé de subimaxe deformada}
    \label{fig:subfigura-deformada}
  \end{subfigure}
  \caption{Pé de imaxe xeral}
  \label{fig:exemplo-subfiguras}
\end{figure}

\clearpage

\section{Inclusión de táboas}

Se precisamos táboas no noso documento, incluirémolas do xeito que se
indica na táboa~\ref{tab:exemplo} (páxina~\pageref{tab:exemplo}). Se
o facemos así, \LaTeX ubicará cada táboa no mellor lugar posible,
lugar que pode variar a medida que o documento vaia crecendo coa
inclusión de máis texto e outros elementos (máis imaxes, táboas,
etc.).

\begin{table}[hp!]
  \centering
  \rowcolors{2}{white}{udcgray!25}
  \begin{tabular}{c|c}
  \rowcolor{udcpink!25}
  \textbf{Título de columna} & \textbf{Outro título de columna} \\\hline
  \textit{Título de fila} & Contido de celda \\
  \textit{Título de fila} & Contido de celda \\
  \textit{Título de fila} & Contido de celda \\
  \textit{Título de fila} & Contido de celda \\
  \textit{Título de fila} & Contido de celda \\
  \textit{Título de fila} & Contido de celda \\
  \end{tabular}
  \caption{Pé de táboa descritivo}
  \label{tab:exemplo}
\end{table}

Para táboas longas que ocupan varias páxinas, como é o caso da \ref{tab:longa},
recoméndase o uso do paquete \texttt{lontable}, descomentando a liña
correspondente no ficheiro raíz do proxecto (\verb+memoria_tfg.tex+).

\rowcolors{2}{white}{udcgray!25}
\begin{longtable}{l|r|c}
  \caption{Pé descritivo dunha táboa longa}
  \label{tab:longa} \\

  \rowcolor{udcpink!25}
  \textbf{Primeira columna} & \textbf{Segunda columna} & \textbf{Terceira columna} \\\hline
  \endfirsthead

  \multicolumn{3}{c}{\tablename\ \thetable{} -- {\small \textit{(vén da páxina anterior)}}} \\
  \rowcolor{udcpink!25}
  \textbf{Primeira columna} & \textbf{Segunda columna} & \textbf{Terceira columna} \\\hline
  \endhead

  \multicolumn{3}{c}{\dotfill{\small \textit{(continúa na páxina seguinte)}}\dotfill} \\
  \endfoot

  \endlastfoot

  Texto de exemplo & abcdef ghjijklmn & 123.456778 \\
  Texto de exemplo & abcdef ghjijklmn & 123.456778 \\
  Texto de exemplo & abcdef ghjijklmn & 123.456778 \\
  Texto de exemplo & abcdef ghjijklmn & 123.456778 \\
  Texto de exemplo & abcdef ghjijklmn & 123.456778 \\
  Texto de exemplo & abcdef ghjijklmn & 123.456778 \\
  Texto de exemplo & abcdef ghjijklmn & 123.456778 \\
  Texto de exemplo & abcdef ghjijklmn & 123.456778 \\
  Texto de exemplo & abcdef ghjijklmn & 123.456778 \\
  Texto de exemplo & abcdef ghjijklmn & 123.456778 \\
  Texto de exemplo & abcdef ghjijklmn & 123.456778 \\
  Texto de exemplo & abcdef ghjijklmn & 123.456778 \\
  Texto de exemplo & abcdef ghjijklmn & 123.456778 \\
  Texto de exemplo & abcdef ghjijklmn & 123.456778 \\
  Texto de exemplo & abcdef ghjijklmn & 123.456778 \\
  Texto de exemplo & abcdef ghjijklmn & 123.456778 \\
  Texto de exemplo & abcdef ghjijklmn & 123.456778 \\
  Texto de exemplo & abcdef ghjijklmn & 123.456778 \\
  Texto de exemplo & abcdef ghjijklmn & 123.456778 \\
  Texto de exemplo & abcdef ghjijklmn & 123.456778 \\
  Texto de exemplo & abcdef ghjijklmn & 123.456778 \\
  Texto de exemplo & abcdef ghjijklmn & 123.456778 \\
  Texto de exemplo & abcdef ghjijklmn & 123.456778 \\
  Texto de exemplo & abcdef ghjijklmn & 123.456778 \\
  Texto de exemplo & abcdef ghjijklmn & 123.456778 \\
  Texto de exemplo & abcdef ghjijklmn & 123.456778 \\
  Texto de exemplo & abcdef ghjijklmn & 123.456778 \\
  Texto de exemplo & abcdef ghjijklmn & 123.456778 \\
  Texto de exemplo & abcdef ghjijklmn & 123.456778 \\
  Texto de exemplo & abcdef ghjijklmn & 123.456778 \\
  Texto de exemplo & abcdef ghjijklmn & 123.456778 \\
  Texto de exemplo & abcdef ghjijklmn & 123.456778 \\
  Texto de exemplo & abcdef ghjijklmn & 123.456778 \\
  Texto de exemplo & abcdef ghjijklmn & 123.456778 \\
  Texto de exemplo & abcdef ghjijklmn & 123.456778 \\
  Texto de exemplo & abcdef ghjijklmn & 123.456778 \\
  Texto de exemplo & abcdef ghjijklmn & 123.456778 \\
  Texto de exemplo & abcdef ghjijklmn & 123.456778 \\
  Texto de exemplo & abcdef ghjijklmn & 123.456778 \\

\end{longtable}


\section{Inclusión de código fonte}

Se precisamos incluír fragmentos de código fonte, podemos facelo, por exemplo, da
seguinte maneira:

\begin{lstlisting}[language=C]
#include <stdio.h>
#define N 10

int main()
{
  int i;

  // Isto é un comentario
  puts("Ola, mundo!");

  for (i = 0; i < N; i++)
  {
    puts("LaTeX é a ferramenta de edición ideal para profesionais da informática!");
  }

  return 0;
}
\end{lstlisting}

\section{Uso da relación de acrónimos e do glosario}

Os acrónimos edítanse no ficheiro \texttt{bibliografia/acronimos.tex}
e úsanse empregando a orde \texttt{acrlong} para obter o termo
completo (deste xeito: \acrlong{erlang}), a orde \texttt{acrshort}
para obter o acrónimo (deste xeito: \acrshort{erlang}). A primeira vez
que usamos un termo con acrónimo no documento é recomendable usar orde
\texttt{acrfull} (que produce ambas versións á vez:
\acrfull{erlang}). Os acrónimos que non se usan no documento, non
aparecen na relación que se xerar na versión PDF.

Pola súa banda, os termos do glosario edítanse no ficheiro
\texttt{bibliografia/glo\-sa\-rio.tex} e úsanse empregando a orde
\texttt{gls} (deste xeito, \gls{bytecode}) ou \texttt{Gls} (deste
xeito, \Gls{bytecode}). Ao igual que os acrónimos, os termos que non
se usan no documento, non aparecen na relación que se xera na versión
PDF.

%\include{contido/...}
 \chapter{Conclusións}
\label{chap:conclusions}

\lettrine{D}{erradeiro} capítulo da memoria, onde se presentará a
situación final do traballo, as leccións aprendidas, a relación coas competencias da titulación en xeral e a mención en particular,
posibles liñas futuras,\dots

\section{Resultados}\label{chap:resultados}
Tras finalizar o proxecto, pódese garantir que engadir novas extensións coas súas correspondentes novas instrucións non comprometen o traballo anterior. O funcionamento do resto de módulos segue sendo correcto, o rendemento non se viu deteriorado en ningún momento.

Por outra parte, a posibilidade de modificar as latencias dalgunhas instrucións grazas á parametrización engadida, mostra como cambia o rendemento no conxunto dun programa. Por exemplo, á hora de executar o benchmark SPMV para que realice multiplicación de enteiros, obtemos resultados moi interesantes segundo as latencias. Como vemos na táboa \ref{tab:rendemento_spmv}, 

\begin{table}[hp!]
    \centering
    \rowcolors{2}{white}{udcgray!25}
    \begin{tabular}{c|c|c}
    \rowcolor{udcpink!25}
    \textbf{Modificacións realizadas} & \textbf{Número de instrucións}  & \textbf{Tempo} 
    \\\hline
    \textit{SPMV base} & a & a \\
    \textit{Latencia de mul = 5} & a & b\\
    \textit{Latencia de mul = 10} & a & b\\
    \textit{Latencia de mulh = 5} & a & b\\
    \textit{Latencia de mul = 5 e mulh = 5} & a & b\\
    \end{tabular}
    \caption{Rendemento do benchmarks SPMV segundo as latencia de distintas operacións.}
    \label{tab:rendemento_spmv}
\end{table}

Finalmente, destacar que a velocidade de simulación deste proxecto en comparación coa que se podería obter se VHDL ou Verilog fose empregado é moi superior.

\section{Traballo futuro}\label{chap:traballo_futuro}
Agora mesmo, o simulador inclúe todas as instrucións implementadas ata o que é conecido como extensión 'G'. Poderíanse engadir máis extensións, así como unha memoria ¿? ou incluir soporte para 64bits.




 %%%%%%%%%%%%%%%%%%%%%%%%%%%%%%%%%%%%%%%%
 % Apéndices, glosarios e bibliografía  %
 %%%%%%%%%%%%%%%%%%%%%%%%%%%%%%%%%%%%%%%%

 \appendix
 \appendixpage
 \appendix
\chapter{Material adicional}
\label{chap:adicional}

\lettrine{E}{ste} capítulo ten formato de apéndice, inclúe material adicional que non ten cabida no corpo principal do documento, como código de tests ou exemplos de módulos.

\section{Exemplo de código de probas}
\label{cod_test}

\begin{lstlisting}[language=C]
#include <stdio.h>
#include <stdint.h>

// Función para emular la operación remu en RISC-V
int32_t remu(int32_t a, int32_t b) {
    int32_t result = 0;
    asm volatile ("remu %0, %1, %2" : "=r"(result) : "r"(a), "r"(b));
    return result;
}

// Funcion de prueba
void TEST_REMU(int id, int32_t expected, int32_t a, int32_t b, int32_t *res) {
    int32_t result = remu(a, b); 
    if (result != expected) {
        (*res)++;
    }
}

int main(){
  int res = 0;

  TEST_REMU( 1,  2,   20,   6, &res );
  TEST_REMU( 2,  2,  -20,   6, &res );
  TEST_REMU( 3,  20,  20,  -6, &res );
  TEST_REMU( 4, -20, -20,  -6, &res );

  TEST_REMU( 5,      0, -1<<31,  1, &res );
  TEST_REMU( 6, -1<<31, -1<<31, -1, &res );

  return res;
}
\end{lstlisting}

\section{Módulo de multiplicación}
\label{mul_module}
A continuación móstrase o código do módulo de multiplicación, composto polo ficheiro de cabeceira e o correspondente corpo. 

\begin{lstlisting}[language=C++]
#ifndef MUL_H
#define MUL_H

#include "systemc.h"
#include "structsRV.h"
#include "config.h"
#include "auxFuncs.h"


SC_MODULE(mul) {
public:

	sc_in <bool> clk, rst;
	sc_in <instruction> I;

	// Hazard detection from Decod
	sc_in <sc_uint<5>> rs1In, rs2In;
	sc_out <bool> hzrdRs1Out, hzrdRs2Out;

	sc_out <bool> readyFenceMulOut;

	sc_out <instruction> instOut;


	void multiplication();

	void hazardDetection();

	SC_CTOR(mul) {
		cout << "mul: " << name() << endl;

		// NOP
		INST = createNOP();

		SC_METHOD(multiplication);
		sensitive << clk.pos();

		SC_METHOD(hazardDetection);
		sensitive << rs1In << rs2In << fire;

		fire.write(true);

	}

private:

	instruction INST;
	sc_signal <bool> fire;

	instruction pipeline[pipelineSizeMul];

	bool pipelineFull = false;
	bool flagDiv = false;

};

#define MUL 16
#define MULH 17
#define MULHSU 18
#define MULHU 19
#define DIV 20
#define DIVU 21
#define REM 22
#define REMU 23

#endif
\end{lstlisting}

\begin{lstlisting}[language=C++]
#include "mul.h"
#include "alu.h"


// COMPLETELY SEGMENTED
void mul::multiplication() { 

	sc_int <32> A, B, res;
	sc_uint <5> opCode;
	short target;
	double tiempo;

	tiempo = sc_time_stamp().to_double() / 1000.0;

	if (rst.read()) {

		// NOP
		INST = createNOP();
		instOut.write(INST); 

		// empty pipeline
		for (int i = 0; i < pipelineSizeMul; i++) {
			pipeline[i] = createNOP();
		}


	} else {
		
		// Get data
		INST = I.read();

		A = INST.opA;
		B = INST.opB;
		target = INST.rd;
		opCode = INST.aluOp;

		// Independant pipeline for each instruction 
		int cyclesInPipeline = 0;
		instruction output = pipeline[0];

		for (int i = 0; i < pipelineSizeMul - 1; i++) {

			cyclesInPipeline = pipelineSizeMul - i;

			if (pipeline[i].wReg && getLatencyOp(pipeline[i].aluOp,pipeline[i].target) <= cyclesInPipeline) {

				output = pipeline[i];
				pipeline[i] = createNOP();

				// Div in output
				if (pipeline[i].aluOp == DIV || pipeline[i].aluOp == DIVU || 
					pipeline[i].aluOp == REM || pipeline[i].aluOp == REMU ) {
					flagDiv = false;
					pipelineFull = false;
				}
				break;
			}
		}

		if (output.aluOp != 0) {
			pipelineFull = false;
		}

		instOut.write(output);

		// Loop to shift pipeline content
		// Pos 0: exit
		// Pos latencyMUL-1: newElement
		for (int i = 0; i < pipelineSizeMul - 1; i++) {
			pipeline[i] = pipeline[i + 1];
		}

		sc_int<64> tmp = 0;

		// Operate
		switch (opCode)
		{
		case MUL: 
			tmp = ((sc_int<32>)A) * ((sc_int<32>)B);
			INST.aluOut = INST.dataOut = tmp(31,0);
			strcpy(INST.desc, "mul");
			break;

		case MULH: 
			tmp = ((sc_int<32>)A) * ((sc_int<32>)B);
			INST.aluOut = INST.dataOut = tmp(63,32);
			strcpy(INST.desc, "mulh");
			break;

		case MULHU:
			tmp = ((sc_uint<32>)A) * ((sc_uint<32>)B);
			INST.aluOut = INST.dataOut = tmp(63, 32);
			strcpy(INST.desc, "mulhu");
			break;

		case MULHSU:
			tmp = ((sc_int<32>)A) * ((sc_uint<32>)B);
			INST.aluOut = INST.dataOut = tmp(63, 32);
			strcpy(INST.desc, "mulhsu");
			break;

		case DIV:
			if (B == 0) {
				cerr << "Divider can't be 0 " << endl;
			}
			INST.aluOut = INST.dataOut = ((sc_int<32>)A) / ((sc_int<32>)B);
			strcpy(INST.desc, "div");
			flagDiv = true;
			break;

		case DIVU:
			if (B == 0) {
				cerr << "Divider can't be 0 " << endl;
			}
			INST.aluOut = INST.dataOut =((sc_uint<32>)A) / ((sc_uint<32>)B);
			strcpy(INST.desc, "divu");
			flagDiv = true;
			break;

		case REM:
			if (B == 0) {
				cerr << "Divider can't be 0 " << endl;
			}
			INST.aluOut = INST.dataOut = ((sc_int<32>)A) % ((sc_int<32>)B);
			strcpy(INST.desc, "rem");
			flagDiv = true;
			break;

		case REMU:
			if (B == 0) {
				cerr << "Divider can't be 0 " << endl;
			}
			INST.aluOut = INST.dataOut =  ((sc_uint<32>)A) % ((sc_uint<32>)B);
			strcpy(INST.desc, "remu");
			flagDiv = true;
			break;

		default:
			INST = createNOP();
			break;
		}

		// New instruction
		pipeline[pipelineSizeMul - 1] = INST;

	} 
	fire.write(!fire.read());
}

void mul::hazardDetection() {

	int rs1 = rs1In.read();
	int rs2 = rs2In.read();


	bool aux1 = false, 
		 aux2 = false;

	int cont = 0;
	bool emptyPipeline = false;

	// Prevents RAW
	for (int i = 0; i < pipelineSizeMul; i++) {

		if (pipeline[i].wReg) {

			if (rs1 == pipeline[i].rd) {
				aux1 = true;
			}

			if (rs2 == pipeline[i].rd) {
				aux2 = true;
			}
		}
		else {
			cont++;
		}
	}

	if (instOut.read().wReg) {

		if (rs1 == instOut.read().rd) {
			aux1 = true;
		}

		if (rs2 == instOut.read().rd) {
			aux2 = true;
		}
	}
	else {
		emptyPipeline = true;
	}

#if SOLO_1OP

	if (I.read().wReg) {
		int opCode = I.read().aluOp;

		if (isMulModuleOp(opCode)) {
			aux1 = aux2 = true;
			pipelineFull = true;
		}
		else if (!pipelineFull) {
			aux1 = aux2 = false;
		}
	}
	else {
		emptyPipeline = emptyPipeline && true;
	}
	
	
#else

	if (I.read().wReg) {
		int opCode = I.read().aluOp;

		if (flagDiv && isMulModuleOp(opCode)) {
			if (opCode == DIV || opCode == DIVU || opCode == REM || opCode == REMU) {
				aux1 = aux2 = true;
			}
			else {
				aux1 = aux2 = false;
			}
		}
	}
	else {
		aux1 = aux2 = false;
		emptyPipeline = emptyPipeline && true;
	}
#endif
	
	if (cont == pipelineSizeMul && emptyPipeline) {
		readyFenceMulOut.write(true);
	}
	else {
		readyFenceMulOut.write(false);
	}

	hzrdRs1Out.write(aux1);
	hzrdRs2Out.write(aux2);
}
\end{lstlisting}

%\include{anexos/...}

 \printglossary[type=\acronymtype,title=\nomeglosarioacronimos]
 \printglossary[title=\nomeglosariotermos]

 \bibliographystyle{IEEEtranN}
 \bibliography{\bibconfig,bibliografia/bibliografia}
 \clearpage
 
\end{document}

%%%%%%%%%%%%%%%%%%%%%%%%%%%%%%%%%%%%%%%%%%%%%%%%%%%%%%%%%%%%%%%%%%%%%%%%%%%%%%%%
